% =============================================================================
% Appendix C — WTS Sample Memory
% =============================================================================
\chapter{WTS Sample Memory}

\epigraph{``A good sample library is worth a thousand oscillators.''}%
{\textit{--- Studio wisdom}}

\section*{Introduction}

The Wavetable Synthesizer's sample data lives entirely in host memory, outside
both the 6502's 64\,KB address space and the 512\,KB XRAM.  Samples are loaded
from industry-standard SF2 (SoundFont~2) files on the host filesystem and
decoded into PCM audio buffers that the WTS voices reference during playback.

This appendix describes how sample memory is organized, loaded, and used.


% ═══════════════════════════════════════════════════════════════════════════════
\section{Loading Soundfonts}
\label{sec:wts-loading}
% ═══════════════════════════════════════════════════════════════════════════════

Soundfont loading is initiated through the FileIO controller using command \$15
(\texttt{SfLoad}).  The process is:

\begin{enumerate}
  \item Write the soundfont filename (without path) into the FioName buffer at
        \addr{B9B0}.
  \item Write command \$15 to the FileIO command register at \addr{B9A0}.
  \item Poll \texttt{WtsSoundfontStatus} at \addr{A183} until loading completes.
\end{enumerate}

The file must be located in the \texttt{{\raise.17ex\hbox{$\scriptstyle\sim$}}/e6502-programs}
directory (the same directory used for BASIC program files).

\begin{lstlisting}
10 REM load a General MIDI soundfont
20 F$ = "GeneralUser_GS.sf2"
30 FOR I = 1 TO LEN(F$)
40   POKE $B9B0+I-1, ASC(MID$(F$,I,1))
50 NEXT I
60 POKE $B9B0+LEN(F$), 0 : REM null terminator
70 POKE $B9A0, $15 : REM SfLoad command
80 REM wait for load to complete
90 IF PEEK($A183) <> 0 THEN 90
\end{lstlisting}

\begin{notebox}
The MML parser can also trigger automatic soundfont loading when an
\texttt{@I\textless n\textgreater} or \texttt{@D\textless n\textgreater} WTS
instrument command is encountered and no soundfont is currently loaded.
\end{notebox}


% ═══════════════════════════════════════════════════════════════════════════════
\section{SF2 Format Overview}
\label{sec:sf2-format}
% ═══════════════════════════════════════════════════════════════════════════════

SF2 (SoundFont~2) is an industry-standard format originally developed by
Creative Labs.  Each file contains:

\begin{itemize}
  \item \textbf{PCM audio samples} --- raw waveform data, typically 16-bit mono
        or stereo at 44\,100\,Hz.
  \item \textbf{Instrument definitions (presets)} --- each preset maps key ranges
        and velocity layers to one or more sample zones.
  \item \textbf{Articulation data} --- volume envelopes, filter settings, loop
        points, and tuning parameters for each zone.
\end{itemize}

The NovaVM's WTS engine reads and interprets all of these at load time, building
internal data structures for real-time playback.


% ═══════════════════════════════════════════════════════════════════════════════
\section{Instrument Mapping}
\label{sec:wts-instruments}
% ═══════════════════════════════════════════════════════════════════════════════

Each preset in the loaded soundfont becomes a WTS instrument, numbered
sequentially starting from~0.  For General MIDI--compatible soundfonts, the
instrument numbers follow the GM program mapping:

\begin{center}
\small
\begin{tabular}{@{}cl@{\qquad}cl@{}}
\toprule
\textbf{Index} & \textbf{GM Instrument} &
\textbf{Index} & \textbf{GM Instrument} \\
\midrule
0   & Acoustic Grand Piano & 24  & Acoustic Guitar (nylon) \\
4   & Electric Piano 1     & 32  & Acoustic Bass \\
8   & Celesta              & 40  & Violin \\
11  & Vibraphone           & 48  & String Ensemble 1 \\
16  & Drawbar Organ        & 56  & Trumpet \\
19  & Church Organ         & 73  & Flute \\
\bottomrule
\end{tabular}
\end{center}

The total number of available instruments can be queried through the WTS
instrument enumeration interface at \addr{A1A0}.

WTS voices select an instrument by writing the instrument index to the voice's
instrument register (offset~+2 from the voice base address).


% ═══════════════════════════════════════════════════════════════════════════════
\section{Memory Usage}
\label{sec:wts-memory}
% ═══════════════════════════════════════════════════════════════════════════════

Sample data is decoded from the SF2 file at load time and stored as uncompressed
PCM buffers in host memory.  Typical memory usage:

\begin{center}
\small
\begin{tabular}{@{}lc@{}}
\toprule
\textbf{Soundfont Type} & \textbf{Approximate Host RAM} \\
\midrule
Small (single instrument)  & 1--5\,MB \\
Medium (partial GM set)    & 10--30\,MB \\
Full GM soundfont          & 30--100\,MB \\
\bottomrule
\end{tabular}
\end{center}

\begin{warningbox}
Sample memory resides entirely in the host process and does not affect the
6502's 64\,KB address space or the 512\,KB XRAM.  However, very large soundfonts
may consume significant host RAM.
\end{warningbox}


% ═══════════════════════════════════════════════════════════════════════════════
\section{Limitations}
\label{sec:wts-limits}
% ═══════════════════════════════════════════════════════════════════════════════

\begin{itemize}
  \item Only one soundfont can be loaded at a time.  Loading a new soundfont
        replaces the previous one entirely, including all instrument definitions
        and sample data.
  \item Soundfont files must be in SF2 format; compressed variants (SF3/SF4) are
        not supported.
  \item Sample data is read-only from the 6502's perspective --- there is no
        mechanism to upload custom PCM data from CPU RAM or XRAM.
  \item The maximum number of simultaneous WTS voices is~8, regardless of how
        many instruments are defined in the soundfont.
\end{itemize}
